\section{An effective exponent budget and its formal arithmetic core}
\label{hbc-section}

\begin{corollary}[An explicit pure-power exponent budget]\label{hbc-budget}
Let $a,b$ be positive coprime integers, $p>7$ prime, and
$F(a,b)=Q^p$ for an integer $Q>1$. Put $H=\max(a,b)$ and
$T=\log13+4\log H$. Then
\[
 p^p<F(a,b)\le13H^4,\qquad p\log p<T.
\]
If $T\ge e^2$, then
\[
 p<\frac{2T}{\log T}.
\]
In particular $p=O(\log H/\log\log H)$ as $H\to\infty$,
with an effective absolute constant.
\end{corollary}
\begin{proof}
Corollary~\ref{ts-pure-support} gives $Q>p$, and
Corollary~\ref{ts-height} gives the polynomial height bound.
Taking logarithms proves the first two assertions.
For the explicit rearrangement, suppose $p\ge2T/\log T$.
The function $x\log(x/2)$ increases for $x\ge2$, and
$2T/\log T\ge e^2$ when $T\ge e^2$.
For $z=\log T\ge2$, one has $2\log z<z$: the difference
$z-2\log z$ is positive at $2$ and has nonnegative derivative
thereafter. Therefore
\[
 p\log(p/2)\ge\frac{2T}{\log T}\log\frac{T}{\log T}
   =2T\left(1-\frac{\log\log T}{\log T}\right)>T,
\]
contrary to $p\log p<T$. The asymptotic form follows since
$T=4\log H+O(1)$. No rearranged bound is needed when $T<e^2$.
\end{proof}

This budget bounds the exponent in terms of the actual point height.
It supplies no height upper bound. In particular, it does not rule out
arbitrarily high seeds with moving prime factors, or treat nonunit
residuals $F=VQ^p$ by the same fixed-level argument.

\paragraph{Ordinary proofs before kernel verification.}
The reviewed local proof uses the following elementary facts. If
$q\ge7$ and $t^2\le4q$, then
\[
 1-q<t<q-1,
 \qquad (q-1)^2-4q=q(q-6)+1>0.
\]
Hence an even such trace cannot be congruent to $1$ or $-1$
modulo $q$. For odd $p,q$ and even $t$, if $p$ divides either
$q+1-t$ or $q+1+t$, then $2p$ divides that positive integer.
The displayed interval places it below $2q$, giving $p<q$.
The decomposition into one divisible factor uses primality in the
separate ordinary argument; it is not hidden in the formal signature.

The polynomial height estimates also have direct certificates:
\begin{align*}
 (a+b)^4-F(a,b)&=ab(a^2+ab+b^2),\\
 13a^4-F(a,b)&=(a-b)(12a^3+9a^2b+4ab^2+b^3).
\end{align*}
For $a,b\ge0$ the first gives $F\le(a+b)^4$; the second gives
$F\le13a^4$ when $b\le a$. Symmetry gives $F\le13H^4$.

\paragraph{Exact scope of the new Lean module.}
The module \texttt{BoundarySupportArithmetic.lean} checks ten
declarations: the strict Hasse interval, a small-multiple lemma,
the exclusion of unit residues for even traces, doubling an odd
divisor of an even number, the trace-factor prime cutoff, the
sum-height identity, the sum-height inequality, the left-coordinate
height inequality, its maximum-coordinate version, and the
pure-power height implication.

Its last signature explicitly assumes $p<Q$ and
$F(a,b)=Q^p$. Its trace cutoff explicitly assumes the disjunction
$p\mid q+1-t$ or $p\mid q+1+t$. These are ordinary modular/local
conclusions supplied outside the module, not additional axioms or
claims of their complete formalization. The new proofs use the
existing actual quartic definition. A fresh scoped Lean~4.32.0
build accepts all ten declarations and recompiles fifty-five
unchanged dependency declarations. The complete axiom inventory is
contained in \texttt{propext}, \texttt{Classical.choice}, and
\texttt{Quot.sound}.

No Lean proof of modularity, the full-Sturm modular-form algorithms,
the local elliptic representations, the analytic logarithm estimate,
or ABC is claimed by that arithmetic build. The exact software
replays, ordinary independent reviews, and kernel verification remain
separate evidence with separately recorded source hashes.
